\section{Largest-piece residual localization}\label{sec:localization}

Choose a largest selected name \(L\), breaking area ties by a fixed name
order, and put
\[
  q=|T|-|L|.
\]
The residual pieces are nonempty and their total area is \(q\), so there are
at most \(q\) of them.

\subsection{A translation cap}

For finite nonempty \(A,B\subset\ZZ^2\), define
\[
  D(A,B)=\{t\in\ZZ^2:A+t\subseteq B\}.
\]

\begin{lemma}[Translation cap]\label{lem:translation}
For finite nonempty \(A,B\subset\ZZ^2\),
\[
 |D(A,B)|\leq \max(0,|B|-|A|+1).
\]
Equivalently, if \(D(A,B)\) is nonempty, then
\[
 |D(A,B)|\leq |B|-|A|+1.
\]
\end{lemma}

\begin{proof}
The empty case is immediate.  Suppose \(D=D(A,B)\) is nonempty.  Then
\(A+D\subseteq B\).  Choose an integer linear functional \(\ell\) injective
on the finite set \(A+D\); avoiding the finitely many forbidden integer
slopes suffices.  It is then injective on each factor as well.  If
\(\ell(A)=\{a_1<\cdots<a_m\}\) and
\(\ell(D)=\{d_1<\cdots<d_n\}\), the chain
\[
 a_1+d_1<\cdots<a_m+d_1<a_m+d_2<\cdots<a_m+d_n
\]
contains \(m+n-1\) distinct sums.  Therefore
\[
 |B|\geq|A+D|=|\ell(A)+\ell(D)|
 \geq |A|+|D|-1,
\]
and rearrangement proves the claim.
\end{proof}

The sumset mechanism is classical; the lemma is included to make the
algorithmic constants explicit.  The maximum in the unconditional statement
is essential when no translation exists.

\subsection{The active carrier}

Freeze one legal pose \(A\) of \(L\) and let
\[
  b=X(T)-X(A).
\]
Since \(A\subseteq T\), the vector \(b\) is nonnegative.  Its row coordinates
sum to \(q\), and its column coordinates also sum to \(q\).

\begin{lemma}[Active carrier]\label{lem:active}
After a pose \(A\) of \(L\) is fixed, every residual placement participating
in a target-margin solution is confined to at most \(q\) active rows, at
most \(q\) active columns, and therefore at most \(q^2\) target cells.
\end{lemma}

\begin{proof}
At most \(q\) row coordinates and at most \(q\) column coordinates of \(b\)
are positive.  If a residual pose \(Q\) occurs in a tuple whose residual
margins sum to \(b\), nonnegativity gives \(X(Q)\leq b\) coordinatewise.
Thus \(Q\) uses only rows and columns where \(b\) is positive.  It is
contained in
\[
  G_A=T\cap(\text{active rows}\times\text{active columns}),
  \qquad |G_A|\leq q^2.
\]
\end{proof}

The proof uses both containment of every pose in \(T\) and the declared
row/column finite-grid margins.  It is not a statement about continuous
tomography.

\begin{corollary}[Pose bounds]\label{cor:pose-bounds}
The largest piece has at most \(h(q+1)\) legal poses.  After one is fixed,
each residual piece has at most \(hq^2\) margin-compatible candidate poses.
\end{corollary}

\begin{proof}
Apply \cref{lem:translation} to each oriented copy of \(L\) in \(T\):
every nonempty translation set has size at most
\(|T|-|L|+1=q+1\).  Sum over at most \(h\) orientations.  Any residual pose
compatible with the residual margin vector after \(A\) is fixed is a nonempty
subset of the carrier \(G_A\).  The same lemma gives at most \(q^2\)
translations per orientation, and hence at most \(hq^2\) candidate poses.
This catalogue may overcount poses that do not extend to a full fiber tuple.
\end{proof}
